% ARCHIVED DRAFT: root independently audited and incorporated 2026-11-19 17:55 UTC.
% Main part_B.tex is authoritative; this file preserves drafting provenance.
\begin{saresult}{B.078}{Proved under stated hypotheses}{Every valuation conductor ideal, without a rank restriction}
Let $V$ be a valuation domain with fraction field $K$ and infinite
residue field $l$. Let $R\subset V$ be a unital subring containing a
nonzero proper $V$-ideal $J$, and suppose $R/J=k$ is a finite field
of cardinality $q$. No coefficient-field section or characteristic
assumption is imposed. For every unital
commutative extension $R\subset S$, the following are equivalent:
\[
\begin{gathered}
 \text{the standard polynomial co-operations exist uniquely};\\
 \theta_2\text{ is injective};\qquad (R:_S J)=R;\\
 \operatorname{Ann}_{S/R}(J)=0;\qquad \Omega(R,S)=R[X].
                                                        \tag{B78.1}
\end{gathered}
\]
Every positive case has all multivariable polynomial tensor isomorphisms.
The valuation may have arbitrary rank and arbitrary cardinality; the ideal
may be nonprincipal and need not have a value-group infimum.
\end{saresult}
\begin{proof}
$R$ is local with maximal ideal $J$, and $R\cap\mathfrak m_V=J$.
Indeed the unital map from the field $R/J$ to $V/\mathfrak m_V$ is
injective. For $r\in R\setminus J$, choose $a\in R$ with $ra=1+j$,
$j\in J$; the inverse of $1+j$ lies in $1+J\subset R$, since $J$
is an ideal of $V$. Thus $r$ is a unit of $R$.
We first record the ideal identity that replaces a real-valued infimum:
\[
                 J^N(J:_K J^N)=J\qquad(N\ge1).           \tag{B78.2}
\]
Here all ideal products and ideal quotients in the proof are taken over
$V$, unless an $R$-module is explicitly specified. This is the valuation
ideal form of Clifford regularity: Bazzoni--Salce, Theorem 3(3), pp.~139--140,
and the complete elementary proof in Halter--Koch, Lemmas 5.1--5.3 and
Theorem 5.4(1), author-PDF pp.~8--9, were opened
\cite{BazzoniSalce1996,HalterKoch2009}. For completeness we give the proof
needed here. Set $E=(J:J)$; it is a valuation overring of $V$, with maximal
ideal $\mathfrak p$. Every overring of a valuation domain is a localization:
an element outside it has its inverse inside it, its maximal ideal is
contained in $V$, and inverting the elements of $V$ outside that ideal
gives the overring. In particular $\mathfrak p\subset V$.
If $J=aE$ is principal over $E$, then $J^2a^{-1}=J$.
Otherwise $J\mathfrak p=J$: if $a\in J\setminus J\mathfrak p$, comparison
of any $b\in J$ with $a$ in the valuation ring $E$ shows $b\in aE$.
Indeed if $b\notin aE$, then $a/b\in\mathfrak p$ and
$a\in b\mathfrak p\subset J\mathfrak p$, a contradiction.
In this nonprincipal case $J(E:J)=\mathfrak p$. The left side is a proper
ideal of $E$ (otherwise $J$ would be invertible and hence principal in
the local valuation ring). Conversely, for $z\in\mathfrak p$ one has
$zJ\subsetneq J$, since equality would put $z^{-1}$ in $(J:J)=E$.
Choose $a\in J\setminus zJ$. Comparison of ideals gives $zJ\subset aE$,
so $z/a\in(E:J)$ and $z\in J(E:J)$. Consequently
$J^2(E:J)=J\mathfrak p=J$. In both cases $J=J^2A$ for a fractional
$V$-ideal $A$. Iteration gives $J=J^N A^{N-1}$ and
$A^{N-1}\subset(J:J^N)$, proving (B78.2); the reverse inclusion is
the definition of the ideal quotient.

If $J^2=J$, then $J$ is generated as an $R$-ideal by its $N$th powers
for every $N$. For $0\ne r\in J=J^{N+1}$ choose, in a finite expression
of $r$ as a sum of products of $N+1$ elements of $J$, a product of
valuation at most $v(r)$. If $y$ is a factor of smallest valuation, then
$(N+1)v(y)\le v(r)$, so $r/y^N\in J\subset R$.
Thus B.069 applies and proves (B78.1) in this case.
We may therefore suppose $J^2\ne J$ and fix
\[
             0\ne c\in J\setminus J^2,\qquad W=V[1/c]=R[1/c].
                                                        \tag{B78.3}
\]
The latter equality follows from $cV\subset J\subset R$.
For every $y\in J\setminus\{0\}$ one has $2v(y)>v(c)$, since the
opposite inequality would give $c\in y^2V\subset J^2$.
Taking a factor of smallest valuation in a product gives
\[
                      J^{2h+1}\subset c^hJ\quad(h\ge0). \tag{B78.4}
\]
All statements about valuations here concern inequalities in the actual
ordered value group, with no Archimedean assumption.

The Vandermonde ideals $D_d(R)$ and the powers of $J$ are cofinal.
For one direction, sufficiently many nodes in the $q$ residue classes
force arbitrarily many differences into $J$, so $D_d\subset J^n$ for
large $d$ and any prescribed $n$. For the other, choose $d+1$ elements
$u_0,\ldots,u_d\in V$ with distinct residues in $l$. The determinant
on the nodes $cu_i\in J$ is $\delta=c^h u$, where
$h=d(d+1)/2$ and $u\in V^\times$. By (B78.4),
$J^{2h+1}\subset c^hJ=\delta J\subset\delta R\subset D_d$.
B.064 therefore identifies the effective coefficient ring as
\[
                         C=\bigcup_{n\ge1}(R:_S J^n).   \tag{B78.5}
\]
If $(R:_S J)=R$, induction gives every term in this union equal to $R$,
and B.064 proves all the positive assertions. It remains to show that
injectivity of $\theta_2$ forces $C=R$.

Assume that injectivity. By B.064 the inclusion $R\to C$ is epic.
Localizing at the nonzero elements of $R$ maps $C$ into $K$: an element
of the $n$th term of (B78.5) has its product with $c^n$ in $R$.
Let $L\subset W$ be its rational image and $T$ its kernel. The map
$R\to L$ is epic, and every element of $T$ is annihilated by some power
of $J$. If all $b\in L\setminus\{0\}$ satisfy $v(b)\ge-v(c)$, then
(B78.4) gives $J^3L\subset J$. The common ideal $J^3L$ is an ideal of
both $R$ and $L$, contains $J^3$, and is proper in $R$. The induced
injective epimorphism $R/J^3L\to L/J^3L$ has a local source with
nilpotent maximal ideal. B.073, or its finite-depth expansion proof,
makes this map an isomorphism; hence $L=R$.
Otherwise choose $b\in L$ with $v(b)<-v(c)$. Then $b^{-1}V\subset J$,
so $V\subset L$, and $c^{-1}=b(cb)^{-1}\in L$ because $(cb)^{-1}\in V$.
It follows that $L=W$. If $L=R$, rationalization is a retraction of
the epimorphism $R\to C$, forcing $C=R$. We may thus assume $L=W$.

Lift $c^{-1}$ to $a\in C$ and set $z=ca$. Since $z-1\in T$, some
$c^n$ annihilates $z-1$, whence $z^n(z-1)=0$. Increasing $n$ to be
positive, $e=z^n$ is an idempotent. It has rational image one and
annihilates $T$: the equality $e=z^N$ holds for every $N\ge n$, and
each element of $T$ is killed by some $c^N$. Thus $eC=W$, canonically;
indeed $c(ea)=e$. The element $1-e\in T$ is killed by some $J^M$, so
$D=(1-e)C$ satisfies $J^MD=0$. The epic map $R\to D$ factors through
$R/J^M$, whose maximal ideal is nilpotent. Applying the same
nilpotent-ideal epimorphism argument to its image gives
\[
                          C=W\times R/I,
             \qquad J^M\subset I\subset R,\quad M\ge1,  \tag{B78.6}
\]
with the diagonal inclusion of $R$.

We need an exact coefficient description and an ideal norm.
Choose a set $B\subset R$ of representatives of $k$, and put
$P(X)=\prod_{b\in B}(X-b)$. Differences between distinct elements of
$B$ are units of $R$, $P(R)\subset J$, and $P'(b)\in R^\times$
for every $b\in B$. The monic basis $X^iP^j$, $0\le i<q$, $j\ge0$, gives
\[
 \Omega(R,K)=
 \bigoplus_{0\le i<q}RX^i\ \oplus\!
 \bigoplus_{\substack{0\le i<q\\j\ge1}}
                (J:_KJ^j)X^iP^j.                         \tag{B78.7}
\]
To prove necessity, first subtract the polynomial of degree below $q$
interpolating the values on $B$; its coefficients are in $R$.
For the resulting $F=\sum_{i<q,\,j\ge1}a_{ij}X^iP^j$ and any
$0\ne y\in J$, let $\delta=\min_{i,j}v(a_{ij}y^j)$.
For $b\in B$, substitute $X=b+yZ$ and divide by a coefficient of
valuation $\delta$. Modulo the maximal ideal of $V$, the resulting
polynomial is
\[
             \sum_{j\ge1}\left(\sum_{i<q}
                  \overline{a_{ij}y^j/\mu}\,\bar b^i\right)
                       (\overline{P'(b)}Z)^j,
                      \qquad v(\mu)=\delta.              \tag{B78.8}
\]
For some $b$ this is nonzero by the invertible Vandermonde matrix on
$k$ and the nonzero residues of $P'(b)$, and it has zero constant term. If $\delta<0$, a residue outside
its finite zero set gives a value of $F$ outside $V$. If $\delta=0$,
its values on the infinite field $l$ would have to lie in the finite
set $\bar\mu^{-1}k$, which is impossible for a nonconstant polynomial.
If $\delta>0$ but $\delta\notin v(J)$, a nonzero residue gives a
positive-valued element of $R$ outside $J$, also impossible.
Thus every $a_{ij}y^j$ lies in $J$. Products of $j$ elements of $J$
are divisible by the $j$th power of a factor of smallest valuation;
hence $a_{ij}J^j\subset J$. This proves necessity; sufficiency is
immediate from $P(R)\subset J$. The coefficient ideals in (B78.7)
all lie in $W$, since $c^j a_{ij}\in J\subset R$.

For a polynomial $F$ vanishing on $B$, define its $J$-norm to be the
$V$-ideal
\[
                   \nu_J(F)=\sum_{i<q,\,j\ge1}a_{ij}J^j.
                                                        \tag{B78.9}
\]
Write $F(b+Y)=\sum_{h\ge1}g_{bh}Y^h$ for $b\in B$.
The same reduction (B78.8), now considered for each fixed $y\in J$,
proves equality of the principal ideals of least coefficient value
after multiplication by the corresponding powers of $y$. Taking the
union over $y\in J$, which computes powers of a valuation ideal,
gives
\[
                         \nu_J(F)=\sum_{b\in B,h\ge1}g_{bh}J^h.
                                                        \tag{B78.10}
\]
This equality uses actual ideals rather than a numerical infimum.

If $I=R$ in (B78.6), then $C=W$ and (B78.7) is an $R$-module direct
sum for $\Omega(R,C)$. Its coefficient module of $P$ is
$U=(J:J)$, with $JU=J$. Choose $\beta\in V$ whose residue is not in
$k$. The tensor $\beta P\otimes P-P\otimes\beta P$ has zero polynomial
image, but its projection to $U\otimes_R U$ is nonzero. Indeed modulo
$J$ it becomes $\bar\beta\otimes1-1\otimes\bar\beta$ in
$(U/J)\otimes_k(U/J)$, which is nonzero since $\bar\beta\notin k$.
This contradicts injectivity.

Finally suppose $I\subsetneq R$ in (B78.6). Define, with the indicated
value conditions,
\[
 H=\{F\in W[X]:F(R)\subset I\},\qquad E_0=R[X]+H,
 \qquad H_0=\{F\in H:F|_B=0\}.
                                                        \tag{B78.11}
\]
The first projection of $\Omega(R,C)$ is exactly $E_0$: lift the
coefficients of its second component to $R[X]$ and subtract.
Set $K_0=\sum_{F\in H_0}\nu_J(F)$, a $V$-ideal. Formula (B78.7)
and the polynomials $aP$, $a\in J^{M-1}$, give
\[
                            J^M\subset K_0\subset J.    \tag{B78.12}
\]
We claim $H\not\subset JE_0$. Suppose the contrary and fix $N\ge1$.
For $F\in H_0$, $b\in B$, and $\lambda\in(J:J^N)$, the substitution
$\psi=b+\lambda P^N$ preserves $R$ and sends $F$ into $H_0$.
All its nonzero $P$-degrees are at least $N$. Moreover, (B78.2), the
fact that products of a valuation ideal are generated by pure powers,
and (B78.10) give
\[
 \sum_{b\in B,\,\lambda\in(J:J^N)}\nu_J(F\circ\psi)
       =\sum_{b,h}g_{bh}(J:J^N)^hJ^{Nh}
       =\sum_{b,h}g_{bh}J^h=\nu_J(F).                    \tag{B78.13}
\]
The coefficients $\lambda$ belong to $W$, since $c^N\lambda\in J$.
By the assumed containment write
$F\circ\psi=A+\sum r_tF_t$ with $A\in JR[X]$, $r_t\in J$,
and $F_t\in H$. Subtracting the interpolants on $B$ from all $F_t$
absorbs only elements of $JR[X]$ into $A$, so we may take $F_t\in H_0$.
Project to $P$-degrees at least $N$. The norm of this part of $A$ is
contained in $J^{N+1}$, while the remaining terms have norm contained
in $JK_0$. Consequently (B78.13) implies
$K_0\subset J^{N+1}+JK_0$. Taking $N=M$ and using (B78.12) gives
$K_0=JK_0$. Iterating and using $K_0\subset J$ gives
$J^M\subset K_0\subset J^{M+1}$, so $J^M=J^{M+1}$.
This forces $J=J^2$: for $r\in J$, the equality
$r^M\in J^M=J^{2M}$ provides a product of $2M$ elements of $J$
of valuation at most $Mv(r)$; its least-valued factor $y$ satisfies
$2v(y)\le v(r)$, hence $r\in y^2V\subset J^2$.
The contradiction proves the claim.

Choose $F\in H\setminus JE_0$. In $\Omega(R,C)$ the two elements
\[
                       f=(F,0),\qquad g=(0,P^M\bmod I)
\]
both survive modulo $J\Omega(R,C)$. For $f$ this follows by first
projection; for $g$ it follows from its monic coefficient and
$I\subset J$. Their pure tensor is therefore nonzero over the residue
field $R/J=k$. But its polynomial image is zero, since the two
factors lie in the distinct product components in (B78.6).
This final contradiction excludes every $C\ne R$.
Since $(R:_S J)\subset C$ and (B78.5) is obtained by iterating this
colon operation, $C=R$ is equivalent to the ideal condition in
(B78.1). The proof and all positive tensor assertions are complete.
\end{proof}
